\section{Related Work}\label{sec:related}

Ordinos~\cite{ordinos2020} introduced verifiable tally-hiding elections.
It keeps the aggregated tally encrypted under threshold Paillier and
computes the result by multi-party computation, so nobody, not even the
talliers, learns intermediate values. Hertel et
al.~\cite{ordinosext2021} extend it to Borda, Hare--Niemeyer,
instant-runoff, and Condorcet methods including Schulze, using
Lipmaa--Toft comparisons. They report $135$ minutes for five
candidates and $9$ days $10$ hours for twenty, with the trustees on
i5-class machines over a local network~\cite[pp.\,6, 12]{ordinosext2021}. Cortier, Gaudry, and Yang~\cite{cgy2022}
give a toolbox over threshold ElGamal with a UC-proven conditional-gate
primitive and support Condorcet--Schulze with ties
(timings in Section~\ref{sec:comparison}).
They note that their adjacency-matrix Condorcet
variants reveal the matrix. Kryvos~\cite{kryvos2022} defines
\emph{publicly} tally-hiding elections, where the talliers learn the tally
and the public learns only the result. It builds them from Groth16 proofs over
Pedersen vector commitments for Borda, majority judgement, instant-runoff,
and Condorcet (the Smith set, not Schulze). A twenty-candidate Condorcet
ballot proof takes under $40$ seconds with a $0.5$\,GB proving
key~\cite[p.\,9]{kryvos2022}, and Condorcet is the result function on
which their tally scales worst~\cite[p.\,8]{kryvos2022}.
PQKryvos~\cite{pqkryvos2026} instantiates the same framework with
lattice-based commitments and hash-based proofs.
Section~\ref{sec:comparison} compares our scheme with these systems.

The batched ballot ties the $m^2$ ciphertexts to one committed vector by
a random linear combination, the batch-verification technique of
Bellare, Garay, and Rabin~\cite{bgr98}, which Bayer and
Groth~\cite{10.1007/978-3-642-29011-4_17} put at the core of their shuffle
argument. Our vector arguments need no trusted setup. Huber et al.~\cite{hklr24} show that general-purpose SNARKs can
prove the validity of an exponential-ElGamal ballot, and Röhr et
al.~\cite{rhk25} cut the cost of the exponentiation gadget that
takes up most of such circuits. Both keep the per-entry re-encryption inside the
circuit.

Several verifiable ranked tallies publish the pairwise counts. VERICONDOR~\cite{vericondor2024} encodes a ballot as a pairwise matrix
with strict preference and indifference, and has no tallying authority
(in the style of DRE-ip). It proves the ballot well formed without a
shuffle, through a permutation proof on row sums ($O(m^2)$ for strict
rankings, $O(m^3)$ with indifference), and it publishes the aggregate
matrix. Its authors discuss
the Italian attack and running Schulze on the published matrix. Haines,
Pattinson, and Tiwari~\cite{haines2019verifiable} give a homomorphic Schulze
tally, machine-checked in Coq over axiomatised primitives. It checks
ballots by shuffling and decrypting them, and it decrypts the pairwise
counts. Yang et
al.~\cite{yang2018ranked} tally ranked-choice ballots over exponential
ElGamal with per-element proofs. Our position-vector validity is simpler than VERICONDOR's
permutation proof and handles ties with $O(m^2)$ proofs of $O(m)$ branches
each. Compared with the VSTTE'19 scheme, we remove both leaks (the
shuffled ballots and the decrypted pairwise counts) and replace the axioms with
the library's theorems.

For single transferable vote, Shuffle-Sum~\cite{shufflesum2009} and
Teague, Ramchen, and Naish~\cite{teague2008stv} address the Italian
attack. Risk-limiting tallies~\cite{jamroga2019rlt} publish a sample of
ballots instead of all of them.

Belenios has machine-checked privacy and verifiability proofs in
EasyCrypt~\cite{cortier2018belenios}. Verified verifiers for
elections~\cite{10.1145/3319535.3354247} and for mix-nets~\cite{9519460,287095}
check election and mix-net outcomes after the fact, and Pattinson
and Tiwari~\cite{10.1007/978-3-319-66107-0_26} give the Schulze method
as evidence-carrying computation. 

Barthe et al.~\cite{5552642} formalise $\Sigma$-protocols and their
composition in CertiCrypt without extraction. Butler et
al.~\cite{butler2021formalising} formalise Schnorr, Chaum--Pedersen, and
Okamoto and binary And/Or compositions in CryptHOL, with zero-knowledge as
equality of distributions, as we do, but without executable code, $n$-ary
compositions, or instantiation. Firsov et al.~\cite{10221929} formalise
$\Sigma$-protocols in EasyCrypt, and Almeida et
al.~\cite{almeida2010certifying} give a certifying compiler whose generated
code is outside the certified boundary. 
